\documentclass[11pt,a4paper]{article}
% Shared preamble for RuPhO English problem statements (match T1 2026 style).
% Use: \input{latex/rupho_en_preamble.tex} after \documentclass.
\usepackage[margin=1in]{geometry}
\usepackage{amsmath,amssymb}
\usepackage{graphicx}
\usepackage{float}
\usepackage{enumitem}
\usepackage{microtype}
\usepackage{caption}
\usepackage{tabularx}
\usepackage{hyperref}

\captionsetup{font=small,labelfont=bf,skip=6pt}

\setlist[itemize]{leftmargin=1.2cm}
\setlist[enumerate]{leftmargin=1.2cm}

% One outer box per subproblem: [id | statement | points] (no inner rules)
\newcommand{\problembox}[3]{%
  \par\addvspace{0.42em}%
  \noindent
  \setlength{\fboxsep}{7.5pt}%
  \setlength{\fboxrule}{0.95pt}%
  \fbox{%
    \begin{minipage}{\dimexpr\linewidth-2\fboxsep-2\fboxrule\relax}
    \begin{tabularx}{\linewidth}{@{}>{\raggedright\arraybackslash\bfseries}p{2.1em} @{\hspace{0.9em}} >{\raggedright\arraybackslash}X @{\hspace{0.9em}} >{\raggedleft\arraybackslash\bfseries}p{2.4em}@{}}
    #1 & #2 & #3
    \end{tabularx}
    \end{minipage}%
  }%
  \par\addvspace{0.32em}%
}

\title{\textbf{Lamp and photodiode}}
\author{\normalsize X24 (2024) --- English translation}
\date{}
\begin{document}
\maketitle

Equipment: Oscilloscope AC generator DC source (2 pcs.) Multimeter Incandescent lamp in holder Resistor $1~Ohm$ Photodiode in a housing connected in series with resistor $1~MOhm$ BNC alligator wires (3 pcs.) Banana alligator wires (6 pcs.) Tripod with two legs Ruler $50~cm$ Black cardboard backing laptop

Attention! You have been given a photodiode connected in series with resistor $1~MOhm$, the resistor is wrapped in electrical tape. It is strictly forbidden to remove the electrical tape and connect the photodiode to the circuit otherwise than in series with this resistor. Failure to comply with this requirement will result in the tour being cancelled. You are prohibited from opening or using anything that is located on your computer outside the folder with your last name. Failure to comply with this requirement will result in the tour being cancelled. The effective (rms) voltage across the lamp should not exceed $11~\text{V}$, otherwise the lamp may burn out. Burnt out lamps are not replaced. The multimeter can only be used in voltmeter and ohmmeter modes.

Fig.~1. Equipment

\begin{figure}[H]
    \centering
    \includegraphics[width=0.72\textwidth]{E2_EN_assets/fig1.jpg}
    \caption{Figure 1.}
\end{figure}

\section*{Instructions for using the program}

The program is used to approximate experimental data with a smoothing function. Use the program only when instructed to do so in the assignment text. Experimental data is written to the files \texttt{V(T)\_sine.txt} or \texttt{V(T)\_stages.txt}. The text of the condition will indicate which file to write the data to. In the first line, the quantities under study and their dimensions are entered in the format ``X, [X] Y, [Y]'', where X, Y are quantities, [X], [Y] are their dimensions. There must be a space after the comma. Next, pairs of values X, Y are entered in each line (X and Y are separated by an arbitrary number of spaces). The decimal separator is a dot.

Fig.~2. Data entry example

\begin{figure}[H]
    \centering
    \includegraphics[width=0.72\textwidth]{E2_EN_assets/fig2.jpg}
    \caption{Figure 2.}
\end{figure}

A few seconds after launch, the program prompts you to select a file to import data. To select a file, enter the number 1 or 2 into the console. After this, the program displays two pictures (see Fig.~3): the first shows the initial points and the graph of the function approximating them, the second shows the graph of the derivative of the approximating function. Note 1: After starting the program, it may take time ($\sim 1~\text{min}$) to process the data. When you hover the cursor over the chart, the cursor coordinates are indicated in the lower right corner of the window. The buttons in the lower left corner of the window are responsible for the following actions: return the graphs to their original form; take a step back; go one step further; chart moving mode (after clicking, drag the chart with the left mouse button); chart zoom mode (after clicking, drag to select a region); display settings (sliders adjust graph size and margins); save as PNG (the saved image is not interactive, so arrange the view before saving). Note 2: Sometimes the buttons do not respond; try other buttons, or close the window and restart the program.

Fig.~3. Data output example

\begin{figure}[H]
    \centering
    \includegraphics[width=0.72\textwidth]{E2_EN_assets/fig3.jpg}
    \caption{Figure 3.}
\end{figure}

Please note that the approximation function is selected for a specific physical dependence described in the condition, therefore, for points that do not correspond to this dependence, the approximation is unacceptable.

\section*{Theoretical information}

Electromagnetic radiation (hereinafter referred to as radiation) is a set of electromagnetic waves of different wavelengths. The energetic luminosity $w$ of a body is the total radiated power per unit surface area.

A gray body is a model in which energetic luminosity and temperature are related by the Stefan--Boltzmann law
\[
w=\varepsilon\sigma T^4,
\]
where $\sigma$ is the Stefan--Boltzmann constant and $0<\varepsilon<1$ is the emissivity. Numerically,
\[
\sigma = 5.67\times 10^{-8}~\mathrm{W}/(\mathrm{m}^2\cdot\mathrm{K}^4).
\]

An incandescent lamp is a light source in which a tungsten filament is heated by an electric current. To limit oxidation, the filament is placed in an evacuated bulb.

The resistance of tungsten depends on temperature; in this problem you may treat this dependence as linear:
\[
R=R_0\bigl(1+\alpha(T-T_0)\bigr),
\]
where $R_0$ is the filament resistance at room temperature $T_0$. For tungsten, take
\[
\alpha = 4.8\times 10^{-3}~\mathrm{K}^{-1}.
\]

When measuring the lamp resistance with an ohmmeter, you may neglect heating relative to room temperature. Take room temperature as fixed, $T_0=295~\mathrm{K}$.

This problem is devoted to the study of the electrical properties of the lamp and photodiode, as well as the emissive and thermal properties of the lamp.

\section*{Part A. Electrical and radiative properties of the lamp (4.3 points)}

\problembox{A1}{Take the current-voltage characteristics of the lamp in the voltage range $[0; 10]~\text{V}$ (at least 25 points). Take at least 15 points in the range $[0; 5]~\text{V}$ Plot a graph of the resulting dependence.}{2.00}

The power supplied to the lamp filament at constant voltage leaves the filament mainly as electromagnetic radiation and (weakly) as heat exchange with the surroundings. Because the bulb is evacuated, heat exchange is small compared to radiation. You may treat the filament as a gray body, so its energetic luminosity satisfies
\[
w=\varepsilon\sigma T^4.
\]

\problembox{A2}{Using the assumptions made at the beginning of this part, construct a linearized graph of the power output of the lamp versus its temperature. Draw conclusions about the correctness of the assumptions made.}{1.00}

\problembox{A3}{Using the graph, determine the absorption capacity of tungsten $\varepsilon$, from which the lamp filament is made. You can assume that the lamp filament has a circular cross-section, the length of its luminous part is $L = (15.5 \pm 0.5)~mm$, and its diameter is $d = (130 \pm 5)~um$.}{0.50}

\problembox{A4}{Plot a graph of the temperature of the lamp filament versus the voltage across it. This graph will be used in Part C.}{0.80}

\section*{Part B. Photodiode study (5.7 points)}

A photodiode is a semiconductor element that is photosensitive. This means that its electrical properties depend on the power of radiation hitting it. A photodiode has two terminals: an anode and a cathode. For the photodiode you were given, the red wire corresponds to the anode, the black wire to the cathode. We will consider the voltage on the photodiode positive if the anode potential is greater than the cathode potential. We will consider the current strength to be positive if the current flows through the photodiode from the anode to the cathode. If the current strength is positive, we will say that the current flows in the forward direction, if it is negative, in the opposite direction. In this part of the problem, we will study how the illumination of a photodiode affects its behavior in an electrical circuit. Attention! The negative side of the DC source is grounded, so it must be connected to the oscilloscope ground.

Fig.~4. Photodiode leads

\begin{figure}[H]
    \centering
    \includegraphics[width=0.72\textwidth]{E2_EN_assets/fig4.jpg}
    \caption{Figure 4.}
\end{figure}

Place the tripod on a black cardboard backing. Its use makes it possible to reduce background illumination of the photodiode. In the future, ignore its influence on the measurements. Fix the photodiode in one leg of the tripod, and in the other leg fix the incandescent lamp under the photodiode (the orientation of the photodiode and the lamp is shown in Fig.~5). The distance between the lamp filament and the photodiode should be approximately $10~cm$.

Fig.~5. Orientation of photodiode and lamp

\begin{figure}[H]
    \centering
    \includegraphics[width=0.72\textwidth]{E2_EN_assets/fig5.jpg}
    \caption{Figure 5.}
\end{figure}

\problembox{B1}{Take the current-voltage characteristics of the photodiode for three different values of radiation power incident on the photodiode, corresponding in your configuration to the lamp voltages $0~\text{V}$, $5~\text{V}$, $7~\text{V}$ (at least 20 points for each value). The current through the photodiode in the forward direction should not exceed $1~\mu\text{A}$. At this point, do not examine photodiode voltages less than $-0.4~\text{V}$. Sketch the measurement setup used. Consider the resistance of the oscilloscope $R_{osc} = 1~MOhm$ to be known. Note. The voltage across a photodiode refers to the voltage across it only, and not the total voltage across it and the resistor connected to it.}{3.60}

Sketch the measurement setup used. Consider the oscilloscope resistance $R_{osc} = 1~MOhm$ known.

\problembox{B2}{Construct graphs of the obtained dependencies on one sheet of graph paper.}{1.90}

\problembox{B3}{Make a conclusion: at what voltages on the photodiode ($U<0$ or $U>0$) can it be used to measure the power of radiation incident on it, almost independently of the voltage applied to it?}{0.20}

\section*{Part C. Thermal properties of the lamp (10.0 points)}

The main goal of this part of the problem is to obtain the dependence of the heat capacity of the lamp filament on its temperature over the widest possible temperature range. According to theoretical information, the total radiation power of a gray body (which we consider the filament of a lamp) depends only on its temperature and size, therefore the radiation power of the lamp completely determines its temperature. Connect the photodiode, oscilloscope and power supply as shown in Fig.~6. Set the voltage on the source $16~\text{V}$ and do not change it in the future. With this connection, the current through the photodiode and, therefore, the voltage on the oscilloscope will depend in some way on the power of the light incident on the photodiode. We will study the dependence of the voltage on the oscilloscope on the temperature of the lamp experimentally (points C8, C10), and then approximate it with some function, which we will use for calculations in the future. Further, the voltage on the oscilloscope channel included in the circuit of Fig.~6 will be denoted by the letter $V$, the voltage on the lamp by $U$ (note that it may differ from the generator readings), time by $t$, and the temperature of the lamp filament by $T$. Orient the lamp and photodiode the same way as in part~B.

Fig.~6. Connection diagram

\begin{figure}[H]
    \centering
    \includegraphics[width=0.72\textwidth]{E2_EN_assets/fig6.jpg}
    \caption{Figure 6.}
\end{figure}

Connect the lamp to the alternating current generator and apply a sinusoidal voltage to it with the highest possible amplitude and frequency $f_1 \approx 20~\text{Hz}$. Make sure that the dependence $V(t)$ has the form of a sinusoid with a frequency $2f_1$ and a constant shift relative to zero. Note: there may be a small constant shift in the default generator output that significantly changes the observed signal on the oscilloscope; remove it using the oscillator's DC LEVEL setting. Now apply a step voltage to the circuit with some constant shift and frequency $f_2 = 40~\text{Hz}$. Verify that the dependence $V(t)$ has a sawtooth shape with frequency $f_2$ and a constant offset.

\problembox{C1}{Qualitatively explain why the received signals $V(t)$ have frequencies $2f_1$ and $f_2$ respectively.}{0.20}

Let us theoretically describe the output signal of an oscilloscope for an arbitrary periodic signal on a lamp $U(t)$ with a frequency $f$. We will assume that the regime has been established, so the temperature change is also periodic. Points C2--C6 allow you to find the dependence $V(t)$ in general form. If you have difficulty deriving the theory, you can refer to C7 and use the expressions there. Let $P(T)$, $R(T)$ be the lamp radiation power and the resistance of the lamp filament, respectively, at a filament temperature $T$. Let us denote the heat capacity of the filament at the same temperature by $C(T)$. We will also assume that at any moment in time the temperatures of all points of the filament are the same.

\problembox{C2}{Write down an equation that determines a small change in the temperature of the lamp filament $dT$ for a short time $dt$ at time $t$, when the lamp temperature is $T$. The equation may contain $C(T)$, $R(T)$, $P(T)$, $U(t)$.}{0.30}

In subsequent paragraphs we will assume that the relative deviation of the lamp filament temperature from the average is small. Suppose that when the signal $U(t)$ is applied to the lamp, the filament temperature makes small fluctuations near a certain value $T_1$. Taking into account the approximation made, the following functions can be considered constants at any time: $P(T) \approx P(T_1) = P_1$, $R(T) \approx R(T_1) = R_1$, $C(T) \approx C(T_1) = C_1$.

\problembox{C3}{Express $P_1$ in terms of $\langle U^2(t)\rangle$ and $R_1$ by integrating the equation of part C2 over one period. Here $\langle U^2(t)\rangle$ denotes the time average of $U^2(t)$ over one period: $\displaystyle \langle U^2(t)\rangle = f\int_0^{1/f} U^2(t)\,dt$.}{0.30}

Let at time $t = 0$ the temperature of the lamp filament be equal to $T_1$, and write $\Delta T = T-T_1$.

\problembox{C4}{Find $\Delta T(t)$. Express your answer using $U(t)$, $\langle U^2(t) \rangle$, $R_1$, $C_1$. The answer may contain definite integrals.}{0.30}

Due to the small temperature fluctuations, the voltage on the oscilloscope $V$ also makes small fluctuations near a certain steady-state value $V_1$. Write $\Delta V = V-V_1$. Let the dependence of the voltage on the oscilloscope on the temperature of the filament $V(T)$ and, accordingly, the dependence of the derivative on temperature $(dV/dT)(T)$ be known (in the future, these dependences will be obtained experimentally).

\problembox{C5}{Express $\Delta V$ at an arbitrary time in terms of $\Delta T$ and $(dV/dT)(T_1)$.}{0.10}

\problembox{C6}{Get an explicit expression for $V(t)$. Consider that $V(0) = V_1$. Express your answer using $U(t)$, $\langle U^2(t) \rangle$, $R_1$, $C_1$, $(dV/dT)(T_1)$. The answer may contain definite integrals.}{0.20}

Now, having obtained the dependence in general form, let us consider two special cases of the type of signal on the lamp $U(t)$.

\problembox{C7}{Obtain $V(t)$ for: (i) sinusoidal drive $U(t)=U_0\sin(2\pi ft)$; (ii) a periodic two-level (``step'') drive where $U(t)$ alternates between constants $U_1$ and $U_2$ on successive half-periods of duration $1/(2f)$, with integer period index $N$ and $|U_1|>|U_2|$. Sketch $V(t)$ in each case and label characteristic parameters. Express the filament heat capacity $C_1$ in each case using $(dV/dT)(T_1)$, $R_1$, $f$, $U_0$, $U_1$, $U_2$, and the oscilloscope voltage swing $\Delta V=V_{\max}-V_{\min}$. If needed, you may use $C_1=\dfrac{U_0^2}{\pi f R_1}\dfrac{(dV/dT)(T_1)}{\Delta V}$ (sinusoidal) or $C_1=\dfrac{U_1^2-U_2^2}{fR_1}\dfrac{(dV/dT)(T_1)}{\Delta V}$ (two-level); state explicitly if you use these shortcuts.}{1.00}

Now for further calculations it is necessary to obtain the experimental dependence $V(T)$, approximate it with some function and, using a computer program, obtain graphs of the function itself and its derivative. From theory it follows that when connecting a photodiode according to the diagram in Fig.~6, as the radiation power incident on the photodiode increases, the voltage $V$ increases, but the maximum possible voltage $V$ is limited to the value $\approx8~\text{V}$. If you continue to increase the radiation power after reaching this value, the voltage $V$ will change very little (you can verify this yourself), so voltages $V<8~\text{V}$ should be used for measurements. In order to subsequently use the graph $V(T)$ as efficiently as possible, the voltage $V$ should be close to the limit $8~\text{V}$ (we will use $7~\text{V}$ with a margin) at the maximum possible radiation power of the lamp in each operating mode.

Connect the lamp to the generator with a sinusoidal signal of maximum amplitude and frequency $f_0 = 1~\text{kHz}$. This corresponds to the maximum emission power of the lamp for sinusoidal drive from your generator. Position the lamp and photodiode so that the steady-state voltage on the oscilloscope connected to the photodiode circuit is $V_c \approx 7~\text{V}$. While performing steps C8 and C9, do not change the relative position of the lamp and photodiode. Then disconnect the lamp from the AC source, connect it to the DC source, and apply a voltage $U_c$ so that the oscilloscope voltage is again $V_c \approx 7~\text{V}$.

\problembox{C8}{Read the dependence of the voltage $V$ on the oscilloscope connected to the photodiode on the voltage on the lamp $U$ in the range $[0; U_c]$ (at least 15 points). Using the graph from point A4, convert $U$ to the temperature of the lamp filament $T$. Enter the points of the obtained dependence $V(T)$ into the file \texttt{V(T)\_sine.txt} according to the instructions for working with the program. Temperature $T$ should be in the first column, voltage $V$ in the second. After the end of the round, the file \texttt{V(T)\_sine.txt} should contain the data on which you built the graph. These points should also appear on the answer sheet.}{0.90}

Run the program ``lamp.exe'' and select input data from the file \texttt{V(T)\_sine.txt}. The program will display a graph of the approximated dependence $V(T)$, as well as a graph of its derivative. Apply a sinusoidal voltage to the lamp of the maximum possible amplitude and frequency $f_1 = 20~\text{Hz}$; if necessary, you can change the frequency slightly.

\problembox{C9}{By changing the amplitude of the voltage across the lamp, for each amplitude value, record the values of the oscilloscope average voltage $V_1$ and the difference $\Delta V = V_{max}-V_{min}$. Using dependence $V(T)$, convert $V_1$ to the lamp temperature. Obtain the dependence of the heat capacity of the lamp on temperature in the maximum possible temperature range (at least 10 points).}{1.70}

Now obtain the dependence $V(T)$ again, since the maximum effective voltages on the lamp for sinusoidal and step drive can differ noticeably. Connect the lamp to the generator with a step signal of maximum amplitude and frequency $f_0 = 1~\text{kHz}$; eliminate any DC offset on the generator. This corresponds to the maximum radiation power for step drive. Position the lamp and photodiode so that the steady-state oscilloscope voltage is $V_c \approx 7~\text{V}$. While performing steps C10 and C11, do not change their relative position. Then disconnect the lamp from the AC source, connect it to the DC source, and apply $U_c$ so that the oscilloscope voltage is again $V_c \approx 7~\text{V}$.

\problembox{C10}{Read the dependence of the voltage $V$ on the oscilloscope connected to the photodiode on the voltage on the lamp $U$ in the range $[0; U_c]$ (at least 15 points). Using the graph from point A4, convert $U$ to the temperature of the lamp filament $T$. Enter the points of the obtained dependence $V(T)$ into the file \texttt{V(T)\_stages.txt} according to the instructions for working with the program. Temperature $T$ should be in the first column, voltage $V$ in the second. After the end of the round, the data on which you built the graph should remain in the file \texttt{V(T)\_stages.txt}. These points should also appear on the answer sheet.}{0.90}

Run the program ``lamp.exe'' and select input data from the file \texttt{V(T)\_stages.txt}. The program will display a graph of the approximated dependence $V(T)$, as well as a graph of its derivative. Apply a step voltage to the lamp of the highest possible amplitude and frequency $f_2 = 40~\text{Hz}$; you can change the frequency if necessary.

\problembox{C11}{As you vary the amplitude of the voltage across the lamp, for each pair of amplitude and offset, record the values $V_1$ and $\Delta V = V_{max}-V_{min}$. Using dependence $V(T)$, convert $V_1$ to the lamp temperature. Obtain the dependence of the heat capacity of the lamp on temperature in the maximum possible temperature range (at least 10 points).}{1.70}

\problembox{C12}{Construct graphs of the dependences $C(T)$ obtained in C9, C11. Calculate $dC/dT$ from both graphs, approximating the dependence of the straight line.}{2.40}

\vspace{1em}
\noindent\textit{Source:} \url{https://pho.rs/p/3710}

\end{document}